\documentclass[instructions.tex]{subfiles}
\begin{document}
 
\chapter{Du bon exemple.}
 
\primo Quand il n’y aurait qu’un seul homme de bien, qu’une femme vertueuse dans une ville, dans une paroisse, tous les autres seraient sans excuse de ne pas leur ressembler. Les bons exemples instruisent sans parler, et nous disent : Pourquoi ne vivez-vous pas comme un tel ? Êtes-vous plus délicat que tant de grands seigneurs qui ont vécu dans la pénitence, pendant que vous vivez dans le crime\footnote{Numquid delicatior es illo senatore.} ? Êtes-vous meilleur que tant de filles de qualité qui ont vécu dans la modestie, pendant que vous vivez dans la dissipation ? Cette pensée convertit autrefois saint Augustin.
 
Si c’est une puissante raison de faire le bien quand on le voit faire aux autres, ce n’est pas une raison de faire le mal parce que les autres le font, ni d’y prendre plaisir parce que les autres en rient. Vous jetteriez-vous dans un abîme, parce que les autres iraient s’y précipiter ? Vous amuseriez-vous à plaisanter, si vous voyiez plonger un poignard dans le cœur de votre père ? N’êtes-vous donc pas bien aveugle de vous précipiter en enfer, parce que les autres s’y précipitent ; de plaisanter et de rire, lorsque les autres affligent Dieu, votre Sauveur et votre Père.
 
L’exemple des gens de bien a souvent plus d’efficacité que les sermons d’un prédicateur. La présence d’un homme qui, dans les compagnies, parle avec sagesse, est bien capable d’apprendre à un imprudent ce qu’il doit dire et ce qu’il doit faire. L’exemple d’une femme patiente est plus efficace pour convertir un mari brutal, que les avis d’un pasteur. Les bons exemples dans une famille sont plus utiles pour sanctifier un enfant, que les réprimandes. Pères et mères, vivez en saints, et vos enfans vivront saintement. Voudriez-vous qu’ils vécurent comme des anges, s’ils ne voient chez vous que les œuvres du démon, s’ils n’entendent dans leur famille que le langage de l’enfer ?
 
On convertit plus d’ames par l’exemple que par les miracles, dit saint Chrysostome. Sans les exemples des gens de bien, presque tous les pécheurs périraient. Malheur à ceux qui n’en profitent pas, ou qui s’en raillent !
 
\secundo Tout le monde est obligé de donner bon exemple, parce que tous sont obligés de contribuer au salut d’autrui. Il est difficile de se sauver quand on ne veut sauver que soi-même. Ceux-là y sont obligés, qui ont des devoirs plus particuliers, comme les pères et mères, les maîtres, les juges et les magistrats, les ecclésiastiques et les religieux. Plus l’état est saint, plus les bons exemples doivent être fréquens. Plus vous êtes au-dessus des autres, plus vos exemples doivent être éclatans.
 
Si vous avez été grand pécheur, plus aussi devez-vous donner de grands exemples de vertu. Ce n’est pas assez d’être vertueux, il faut qu’on sache que vous l’êtes, et que, sans affectation de votre part, sans chercher l’estime des autres, l’on voie et l’on sache vos bonnes œuvres\footnote{Videant opera vestra bona, Matth. 5. 16.}. Il ne suffit donc pas à un pécheur d’être converti devant Dieu, il faut qu’il le paraisse devant les hommes. Il n’a pas rougi du vice : malheur à lui, s’il rougit de la vertu !
 
\section{Résolutions}
 
\primo Il faut que, dès ce jour même, je prêche par mes bons exemples.
 
\secundo En conséquence, je m’abstiendrai désormais d’applaudir au mal, de le commettre moi-même par mes discours, mes actions, toute ma conduite ; je remplirai partout mes devoirs avec simplicité, avec courage, avec constance.

\prayer{Mon Dieu, aidez-moi à si bien vivre, que tout en moi répare le passé, et porte mon prochain au bien.}
 
\end{document}
